Here, we provide additional background on the multi-armed bandit problem, and on the baseline algorithms used in this paper. Deeper discussion can be found in \citet{Bubeck-survey12,slivkins-MABbook,LS19bandit-book}.

The \ucb algorithm \citep{bandits-ucb1} explores by assigning each arm $a$ an \emph{index}, defined as the average reward from the arm so far plus a \emph{bonus} of the form $\sqrt{C/n_a}$, where $C = \Theta(\log T)$ and $n_a$ is the number of samples from the arm so far. In each round, it chooses an arm with the largest index. The bonus implements the principle of \emph{optimism under uncertainty}. We use a version of \ucb that sets $C=1$ (a heuristic), which has been observed to have a favorable empirical performance \citep[\eg][]{ZoomingRBA-icml10,RepeatedPA-ec14}.

Thompson Sampling~\citep[][for a survey]{Thompson-1933,TS-survey-FTML18} proceeds as if the arms' mean rewards were initially drawn from some Bayesian prior. In each round, it computes a Bayesian posterior given the history so far, draws a sample from the posterior, and chooses an arm with largest mean reward according to this sample (\ie assuming the sample were the ground truth).
In our setting, the prior is essentially a parameter to the algorithm. We choose the prior that draws the mean reward of each arm independently and uniformly at random from the $[0,1]$ interval. This is one standard choice, achieving near-optimal regret bounds, as well as good empirical performance \citep{Kaufmann-alt12,Shipra-colt12,Shipra-aistats13-JACM}. Each arm is updated independently as a Beta-Bernoulli conjugate prior. Further optimizing \ucb and Thompson Sampling is non-essential to this paper, as they already perform quite well in our experiments.

Provable guarantees for bandit algorithms are commonly expressed via \emph{regret}: the difference in expected total reward of the best arm and the algorithm. Both baselines achieve regret $O(\sqrt{KT\log T})$, which is nearly minimax optimal as a function of $T$ and $K$. They also achieve a nearly instance-optimal regret rate, which scales as $O\rbr{\nicefrac{K}{\Delta}\,\log T}$ for the instances we consider.

The \EpsGreedy algorithm (\pref{fn:eps-greedy}) is fundamentally inefficient in that it does not adaptively steer its exploration toward better-performing arms. Accordingly, its regret rate scales as $T^{2/3}$ (for an optimal setting of $\eps\sim T^{-1/3}$). Fixing such $\eps$, regret does not improve for easier instances.

The \greedy algorithm (\pref{fn:greedy}) does not explore at all, which causes suffix failures. This is obvious when the algorithm is initialized with a single sample ($n=1$) of each arm: a suffix failure happens when the good arm returns $0$, and one of the other arms returns $1$. However, suffix failures are not an artifact of small $n$: they can happen for any $n$, with probability that scales as $\Omega(1/\sqrt{n})$~\citep{BSL-myopic23}.






